\begin{abstract}
Protective perturbations are a promising proactive defense against unauthorized voice cloning, yet diffusion-based purification attacks can remove these perturbations before cloning, bypassing existing defenses. This gap arises because most protection methods optimize against the cloning model directly, treating purification only as an evaluation condition. We propose \method, a purification-aware voice protection framework that incorporates the attacker's diffusion purifier into the optimization loop. \method jointly maximizes post-purification speaker-identity disruption and trajectory deviation inside the reverse diffusion process, making protected waveforms resistant to recovery even after denoising. We instantiate the framework with a De-AntiFake DiffWave purifier and a differentiable ensemble of speaker encoders, and define a benchmark protocol covering direct protection, post-purification robustness, clean-purification controls, and perceptual quality. Our formulation reframes robust voice anti-cloning as an adaptive problem against the attacker's purification path, rather than a direct attack on the cloning backend alone.
\end{abstract}
